\begin{abstract}
A wireless sensor network's lifetime is decided by three choices: which nodes may become cluster heads, which of them serve and who joins whom, and how heads reach the sink. We propose QHP-Net, which gives each choice its own bio-inspired mechanism, each aimed at a failure mode measured in our earlier nine-protocol study. Bacterial quorum sensing with a refractory term picks rested, locally strong candidates. Hippopotamus Optimisation selects heads by minimising the predicted bottleneck stress $\max_i \text{drain}_i/E_i$, and members join whichever head, or the sink directly, minimises their stress. An ant-pheromone link memory learns per-link delivery probability from ARQ acknowledgements alone and drives ETX-aware forwarding. We evaluated QHP-Net against eight protocols, including NSGA-II, the best full-yield protocol of the earlier study, and our earlier HIECF framework, in one engine with energy conservation asserted, over 30 paired topologies. On the lossy channel QHP-Net's first node dies at round 2,133, 249 rounds (13.2\%) after NSGA-II's and 43--52\% after the fuzzy, DQN and HIECF selectors. It wins in all 30 topologies (paired $t(29)=46.2$, $p_{Holm}<10^{-27}$, $d_z=8.4$, confirmed by Wilcoxon), delivers 98.3\% of readings at 16\% less energy per reading, and reaches \fndfrac\% of a physical upper bound. Its gain persists without packet loss (+168 rounds), under heterogeneous energy, with uplink costs charged, and across a nine-cell size sweep. An ablation shows that each selection and clustering mechanism contributes 42--328 rounds, that learned link quality is as good as the true link quality, and that a Physarum congestion router performs worse.
\end{abstract}

\begin{IEEEkeywords}
Wireless sensor networks, network lifetime, green networking, cluster-head selection, quorum sensing, hippopotamus optimization, ant colony, ETX routing, paired t-test, ablation study.
\end{IEEEkeywords}
